% Chapter2/chap2_direction1b.tex
% Encoding: UTF-8

\subsubsection{Map representations for simulation and the role of OpenDRIVE}
A central prerequisite for meaningful simulator-based testing is a map representation that captures not only road geometry but also the structural constraints required for routing, rule interpretation, and reproducible scenario execution. In automated driving research, high-definition (HD) map formats typically encode (i) metric geometry, (ii) lane topology (connectivity and boundaries), and (iii) semantic attributes such as road class, speed limits, and intersection structure. Among widely used representations, \textit{ASAM OpenDRIVE} has become a de-facto standard exchange format for static road networks in driving simulation. OpenDRIVE provides a structured description of roads, lanes, and junctions, and explicitly targets simulation use cases rather than dynamic traffic content \cite{asam_opendrive_170}.

The practical implication for this thesis is that \emph{map validity is not merely visual}: even small structural inconsistencies (e.g., dangling lane successors, inconsistent junction connectivity, implausible curvature transitions) can break routability assumptions or trigger simulator instability. Consequently, rigorous evaluation must treat the map as a \emph{structured graph with geometric constraints}, not as a rendered scene.

\subsubsection{OpenStreetMap as a scalable source and its limitations for lane-level fidelity}
OpenStreetMap (OSM) is attractive for large-scale synthetic testing because it offers broad geographic coverage, continuous community updates, and free accessibility. Prior work has shown that OSM road networks can reach high completeness in many regions, making it suitable for generating street-network graphs at city scale \cite{barringtonleigh2017osm}. At the same time, OSM is a volunteered geographic information (VGI) system, and its quality varies spatially and semantically. Empirical studies emphasize that OSM data quality depends on local contributor activity, resulting in heterogeneous completeness and attribute correctness \cite{haklay2010osm,barron2014osm}.

For simulator maps, the critical gap is that OSM frequently lacks explicit lane-level information and consistent intersection semantics. Even when lane-related tags exist, they may be incomplete, ambiguous, or not aligned with the structural requirements of formats like OpenDRIVE. As a result, OSM$\rightarrow$OpenDRIVE conversion pipelines must introduce heuristics and repair logic to infer missing semantics, resolve junction topology, and enforce connectivity constraints. This makes \emph{the conversion process itself} a major source of domain shift between automatically generated maps and manually authored reference maps.

\subsubsection{Automated HD-map generation and verification: from conversion to auditable quality assurance}
Recent research increasingly treats HD-map generation as an engineering process that must be accompanied by verification and validation (V\&V). Chiang et al.\ propose a systematic V\&V procedure for HD-map generation that separates mapping, metric evaluation, and validation checks, explicitly highlighting the need for repeatable quality criteria \cite{chiang2022hdmapvv}. Complementary work emphasizes open-source pipelines that couple automated generation with rule-based verification to detect structural inconsistencies early (e.g., lane connectivity constraints, invalid geometry, implausible curvature or intersection configurations) \cite{opendrive_lanelet2_vv_2023}.

A key methodological lesson from this literature is that simulator-based testing alone is insufficient as a primary validity oracle. Simulation engines can fail due to resource constraints, runtime assumptions, or nondeterministic behavior under repeated load/unload cycles. Therefore, robust pipelines increasingly rely on \emph{offline-first} validation: structural invariants, deterministic repair logs, and auditable artifacts that can be reproduced independent of the simulator. This aligns directly with the approach in this thesis, where structural and semantic correctness are assessed via deterministic pipeline artifacts, while CARLA-based checks are treated as diagnostic rather than authoritative evidence.

\subsubsection{Implications for domain-gap quantification between map domains}
Because OSM-derived maps require inference and repair, the resulting discrepancy to a manual OpenDRIVE reference is best described as a \emph{structural and semantic domain gap}. Prior OSM literature motivates why such a gap is expected: attribute completeness and correctness vary, and lane-level semantics are not guaranteed \cite{haklay2010osm,barron2014osm}. HD-map V\&V work motivates how to measure and report these differences systematically (e.g., geometry deviations, topology errors, semantic inconsistencies) \cite{chiang2022hdmapvv,opendrive_lanelet2_vv_2023}. In addition, street-network tooling such as OSMnx demonstrates practical methods for extracting and analyzing road graphs from OSM at scale, supporting reproducible preprocessing and network-level statistics \cite{boeing2017osmnx}.

Taken together, related work supports three thesis-level requirements: (i) \emph{deterministic conversion and repair} to make results reproducible, (ii) \emph{explicit structural invariants and quality gates} to separate map validity from simulator behavior, and (iii) \emph{localized analysis} (e.g., tile-based diagnostics) to identify where and why the automatically generated map diverges from the manual reference.
